\section*{ЗАКЛЮЧЕНИЕ}
\addcontentsline{toc}{section}{ЗАКЛЮЧЕНИЕ}

В ходе выполнения курсового проекта была успешно достигнута поставленная цель --- спроектировано и разработано полноценное клиент-серверное веб-приложение для управления проектами и отслеживания задач.

В результате выполнения работы были решены следующие задачи:
\begin{enumerate}
    \item Проведен анализ предметной области, который подтвердил актуальность создания легковесных и интуитивно понятных систем управления задачами (Task Trackers) для малых рабочих групп.
    \item Спроектирована архитектура приложения с использованием микрофреймворка Flask (Python) и объектно-реляционного отображения SQLAlchemy, обеспечивающая высокую производительность и безопасность взаимодействия с базой данных PostgreSQL.
    \item Разработаны пользовательские интерфейсы (Wireframes), которые впоследствии были успешно перенесены в интерактивные HTML-шаблоны с применением CSS-фреймворка Bootstrap 5.
    \item Внедрена надежная подсистема аутентификации и реализована ролевая модель (RBAC), разделяющая функционал системы между «Администраторами» и обычными «Пользователями».
    \item Реализован полный цикл CRUD-операций для трех ключевых сущностей предметной области: Пользователи, Проекты и Задачи.
    \item Интегрирован функционал безопасной загрузки, хранения и скачивания файлов-отчетов, прикрепляемых к конкретным задачам.
    \item Разработан аналитический модуль агрегирования данных, который в реальном времени визуализирует статистику выполнения задач с использованием графической библиотеки Chart.js.
\end{enumerate}

Разработанная система полностью удовлетворяет техническим требованиям, отличается высокой скоростью отклика, защищенностью (шифрование паролей, защита от SQL-инъекций) и адаптивным дизайном, и может служить надежным инструментом для автоматизации рабочих процессов в малых и средних командах.

При разработке пояснительной записки к курсовому проекту строго соблюдались требования \cite{gost_report}, определяющие структуру и правила оформления научно-исследовательских работ.